% Options for packages loaded elsewhere
\PassOptionsToPackage{unicode}{hyperref}
\PassOptionsToPackage{hyphens}{url}
\documentclass[
  ignorenonframetext,
]{beamer}
\newif\ifbibliography
\usepackage{pgfpages}
\setbeamertemplate{caption}[numbered]
\setbeamertemplate{caption label separator}{: }
\setbeamercolor{caption name}{fg=normal text.fg}
\beamertemplatenavigationsymbolsempty
% remove section numbering
\setbeamertemplate{part page}{
  \centering
  \begin{beamercolorbox}[sep=16pt,center]{part title}
    \usebeamerfont{part title}\insertpart\par
  \end{beamercolorbox}
}
\setbeamertemplate{section page}{
  \centering
  \begin{beamercolorbox}[sep=12pt,center]{section title}
    \usebeamerfont{section title}\insertsection\par
  \end{beamercolorbox}
}
\setbeamertemplate{subsection page}{
  \centering
  \begin{beamercolorbox}[sep=8pt,center]{subsection title}
    \usebeamerfont{subsection title}\insertsubsection\par
  \end{beamercolorbox}
}
% Prevent slide breaks in the middle of a paragraph
\widowpenalties 1 10000
\raggedbottom
\AtBeginPart{
  \frame{\partpage}
}
\AtBeginSection{
  \ifbibliography
  \else
    \frame{\sectionpage}
  \fi
}
\AtBeginSubsection{
  \frame{\subsectionpage}
}
\usepackage{iftex}
\ifPDFTeX
  \usepackage[T1]{fontenc}
  \usepackage[utf8]{inputenc}
  \usepackage{textcomp} % provide euro and other symbols
\else % if luatex or xetex
  \usepackage{unicode-math} % this also loads fontspec
  \defaultfontfeatures{Scale=MatchLowercase}
  \defaultfontfeatures[\rmfamily]{Ligatures=TeX,Scale=1}
\fi
\usepackage{lmodern}
\ifPDFTeX\else
  % xetex/luatex font selection
\fi
% Use upquote if available, for straight quotes in verbatim environments
\IfFileExists{upquote.sty}{\usepackage{upquote}}{}
\IfFileExists{microtype.sty}{% use microtype if available
  \usepackage[]{microtype}
  \UseMicrotypeSet[protrusion]{basicmath} % disable protrusion for tt fonts
}{}
\makeatletter
\@ifundefined{KOMAClassName}{% if non-KOMA class
  \IfFileExists{parskip.sty}{%
    \usepackage{parskip}
  }{% else
    \setlength{\parindent}{0pt}
    \setlength{\parskip}{6pt plus 2pt minus 1pt}}
}{% if KOMA class
  \KOMAoptions{parskip=half}}
\makeatother
\setlength{\emergencystretch}{3em} % prevent overfull lines
\providecommand{\tightlist}{%
  \setlength{\itemsep}{0pt}\setlength{\parskip}{0pt}}
\usepackage{bookmark}
\IfFileExists{xurl.sty}{\usepackage{xurl}}{} % add URL line breaks if available
\urlstyle{same}
\hypersetup{
  hidelinks,
  pdfcreator={LaTeX via pandoc}}

\author{\texorpdfstring{}{}}
\date{}

\begin{document}

\begin{frame}[fragile]{The Eight-Stage Pipeline}
\protect\phantomsection\label{the-eight-stage-pipeline}
The build pipeline is orchestrated by
\texttt{scripts/execute\_pipeline.py}, which invokes numbered stage
scripts sequentially. Each stage is a standalone Python script that
exits cleanly or raises an exception to halt the pipeline.

\begin{block}{Stage 00: Environment Setup
(\texttt{00\_setup\_environment.py})}
\protect\phantomsection\label{stage-00-environment-setup-00_setup_environment.py}
Validates the Python environment, checks dependency availability,
creates output directories, and initializes logging. Ensures
\texttt{PYTHONPATH} includes both the repository root and the active
project's \texttt{src/} directory.
\end{block}

\begin{block}{Stage 01: Test Execution (\texttt{01\_run\_tests.py})}
\protect\phantomsection\label{stage-01-test-execution-01_run_tests.py}
Executes pytest with coverage measurement against both infrastructure
tests (\texttt{tests/}) and project tests
(\texttt{projects/\textless{}name\textgreater{}/tests/}). Enforces
configurable failure tolerances:

\begin{itemize}
\tightlist
\item
  \texttt{max\_infra\_test\_failures}: Maximum permitted infrastructure
  test failures (typically 3).
\item
  \texttt{max\_project\_test\_failures}: Maximum permitted project test
  failures (typically 0).
\item
  Coverage thresholds: 60\% infrastructure, 90\% project.
\end{itemize}

The stage generates coverage JSON files for downstream reporting and
saves test results in both JSON and Markdown formats. The infrastructure
test suite alone contains nearly 3,000 tests across 160+ test files.
\end{block}

\begin{block}{Stage 02: Analysis Execution
(\texttt{02\_run\_analysis.py})}
\protect\phantomsection\label{stage-02-analysis-execution-02_run_analysis.py}
Discovers and executes all Python scripts in
\texttt{projects/\textless{}name\textgreater{}/scripts/} in alphabetical
order. Each script is expected to generate figures in
\texttt{output/figures/} and data in \texttt{output/data/}. Scripts
follow the Thin Orchestrator pattern, importing logic from \texttt{src/}
modules. For example, the \texttt{cognitive\_case\_diagrams} project
generates 25+ programmatic figures via 17 DisCoPy renderers during this
stage.
\end{block}

\begin{block}{Stage 03: PDF Rendering (\texttt{03\_render\_pdf.py})}
\protect\phantomsection\label{stage-03-pdf-rendering-03_render_pdf.py}
Compiles Markdown manuscript chapters into a unified PDF via a
three-phase rendering process:

\begin{enumerate}
\tightlist
\item
  \textbf{Pandoc Markdown→LaTeX}: Converts each \texttt{manuscript/*.md}
  file into LaTeX, injecting metadata from \texttt{config.yaml} (title,
  authors, affiliations, DOI).
\item
  \textbf{XeLaTeX Compilation}: Runs \texttt{xelatex} with
  \texttt{biber} for bibliography processing. Handles the
  \texttt{aux}→\texttt{bbl}→\texttt{aux} cycle automatically, with
  cleanup of stale auxiliary files to prevent corruption.
\item
  \textbf{Post-processing}: Applies font embedding verification and
  PDF/A compliance checks.
\end{enumerate}
\end{block}

\begin{block}{Stage 04: Output Validation
(\texttt{04\_validate\_output.py})}
\protect\phantomsection\label{stage-04-output-validation-04_validate_output.py}
Validates the structural integrity of all generated artifacts:

\begin{itemize}
\tightlist
\item
  PDF cross-reference table and trailer verification.
\item
  Figure file existence and minimum size checking.
\item
  Manifest generation with SHA-256 hashes for all output files.
\item
  Markdown structural validation (heading hierarchy, link integrity).
\end{itemize}
\end{block}

\begin{block}{Stage 05: Output Organization
(\texttt{05\_copy\_outputs.py})}
\protect\phantomsection\label{stage-05-output-organization-05_copy_outputs.py}
Copies finalized artifacts to standardized output locations and
generates the pipeline completion manifest.
\end{block}

\begin{block}{Stage 06: LLM Review (\texttt{06\_llm\_review.py})}
\protect\phantomsection\label{stage-06-llm-review-06_llm_review.py}
Invokes a local LLM (via Ollama) to generate:

\begin{itemize}
\tightlist
\item
  \textbf{Executive summary}: A 1-page high-level overview of the
  manuscript.
\item
  \textbf{Quality review}: Detailed feedback on structure, citations,
  and argumentation.
\item
  \textbf{Translations} (optional): Machine translations of the abstract
  into configured target languages.
\end{itemize}

This stage is skippable via configuration and gracefully handles Ollama
unavailability.
\end{block}

\begin{block}{Stage 07: Executive Report
(\texttt{07\_generate\_executive\_report.py})}
\protect\phantomsection\label{stage-07-executive-report-07_generate_executive_report.py}
Aggregates all pipeline metrics---test results, coverage percentages,
rendering duration, validation status, LLM review scores---into a
comprehensive executive report in both JSON and Markdown formats.
\end{block}
\end{frame}

\begin{frame}[fragile]{The Interactive Orchestrators}
\protect\phantomsection\label{the-interactive-orchestrators}
\begin{block}{\texttt{run.sh}: The Standard Orchestrator}
\protect\phantomsection\label{run.sh-the-standard-orchestrator}
The primary user interface is \texttt{run.sh}, a Bash TUI (text user
interface) that presents an interactive menu for pipeline execution.
Features include:

\begin{itemize}
\tightlist
\item
  \textbf{Project selection}: Execute a single project or all discovered
  projects.
\item
  \textbf{Mode selection}: Fast (skip infra tests + LLM), Core (skip LLM
  only), Full (all stages).
\item
  \textbf{Non-interactive mode}:
  \texttt{./run.sh\ -\/-pipeline\ -\/-project\ all\ -\/-core-only} for
  CI/CD integration.
\item
  \textbf{Real-time progress}: Stage timing and status indicators.
\end{itemize}
\end{block}

\begin{block}{\texttt{secure\_run.sh}: The Steganographic Superset}
\protect\phantomsection\label{secure_run.sh-the-steganographic-superset}
The \texttt{secure\_run.sh} orchestrator is a strict superset of
\texttt{run.sh}: it executes the standard eight-stage pipeline and then
appends steganographic post-processing. For each rendered PDF, it
applies metadata injection, cryptographic hashing, alpha-channel text
overlay, and QR code injection, producing a provenance-embedded output
alongside the original. It supports the same project selection and mode
flags as \texttt{run.sh}.
\end{block}
\end{frame}

\end{document}
